\documentclass[10pt,a4paper]{amsart}
\usepackage[margin=19mm]{geometry}
\usepackage{amsmath,amssymb,mathtools}
\usepackage[scr=rsfs,cal=euler]{mathalfa}
\usepackage{xfrac}
\usepackage[unicode,hidelinks,bookmarks=false]{hyperref}
\newcommand{\ie}{i.e.\ }
\newcommand{\cf}{cf.\ }
\newcommand{\resp}{respectively\ }
\newcommand{\Subsection}{\subsection}
\providecommand{\nobreakdash}{\nobreak\hbox{-}}
\setlength{\emergencystretch}{2em}
\multlinegap0pt
\usepackage{amssymb}
\newtheorem{lemma}{Lemma}
\title{Shortest self-orthogonal and LCD embeddings of linear codes over $\mathbb{F}_q+u\mathbb{F}_q$}
\author{Junmin An, Jon-Lark Kim}
\date{Source: arXiv:2608.28222v1 (CC BY 4.0)}
\begin{document}
\maketitle

\noindent\emph{Source and licence.} Shortest self-orthogonal and LCD embeddings of linear codes over $\mathbb{F}_q+u\mathbb{F}_q$. Version arXiv:2608.28222v1 (CC BY 4.0), distributed by arXiv under the Creative Commons Attribution 4.0 International licence, http://creativecommons.org/licenses/by/4.0/, as recorded in the arXiv licence metadata for this exact version and shown on the abstract page https://arxiv.org/abs/2608.28222v1. Full licence notice: \texttt{licenses/arxiv-2608.28222-CC-BY-4.0.txt}.\par\smallskip

\noindent\textit{Editorial scope.} Counted unit: the article's lemma nuzerodiag (source0.tex lines 330--338). The article's complete original proof of it is delivered with it (source0.tex lines 339--419, one \texttt{proof} environment of 81 lines). No mathematics is written, repaired or re-derived: the statement and the proof are the article's own lines, in the article's own notation, and no retained line is substituted or re-flowed.  No further source-local context is needed: every cross-reference of every retained line resolves inside the two delivered spans.  Nothing was omitted from any retained span, so the package carries no omission record.  No retained line contains a citation, so the wrapper carries no bibliography and no bibliography entry.  No retained line contains a figure environment, an image-inclusion command or drawing code, so no image is delivered and the package declares no asset dependency.  Editorial formatting only, in addition: an AMS \texttt{amsart} preamble that restates the article's own declarations and macros used by the retained lines, namely the environments \texttt{lemma} on separate counters, together with the standard packages those lines need.  The retained environments print in this document as Lemma 1 in order of appearance, whereas the article prints the same items with the numbering of its own full document.  The complete native source of the article as distributed in the arXiv e-print for this exact version is bundled unchanged as \texttt{sources/208770/source0.tex}.
\begin{lemma}\label{nuzerodiag}
    Let $R=\mathbb{F}_q+u\mathbb{F}_q$ with $q$ even. For even integers $m_1> 0$ and $m_2\ge 0$,
    \[
    \nu(J_{m_1}\oplus uJ_{m_2})=\begin{cases}
        m_1+m_2, & \mbox{if }m_2>0,\\
        m_1+1, & \mbox{if }m_2=0.
    \end{cases}
    \]
\end{lemma}

\begin{proof}
We first show that
\[
\nu(J_{m_1}\oplus uJ_{m_2})\le\begin{cases}
        m_1+m_2, & \mbox{if }m_2>0,\\
        m_1+1, & \mbox{if }m_2=0.
    \end{cases}
\]
Let $m_2=0$ and $t=m_1/2$. Consider the even weight code $E_{m_1+1}$ over $\{0, 1\}\subseteq \mathbb{F}_q$ of length $m_1+1$. Since $m_1+1$ is an odd integer, $E_{m_1+1}$ does not contain the all-one vector $\mathbf{1}$. Thus,
\[
\mathrm{Hull}(E_{m_1+1})=E_{m_1+1}\cap\langle \mathbf{1}\rangle=\{0\}.
\]
It follows that the standard inner product on $E_{m_1+1}$ is a nondegenerate alternate bilinear form. Let $e_1,\, f_1, \ldots, e_s,\, f_s$ be the symplectic basis of $E_{m_1+1}$ where $s=m_1/2$. Let $A$ be the matrix whose rows are $e_1,\, f_1, \ldots, e_s,\, f_s$ in order. Viewing $A$ as a matrix over $R$, we have $AA^T=J_{m_1}$. Thus, $\nu(J_{m_1})\le m_1+1$.

Assume that $m_2>0$. Define
\[
B=\begin{pmatrix}
    A & \mathbf{0}^T\\u\mathbf{1} & 0\\\mathbf{1} & 1
\end{pmatrix}
\]
where $\mathbf{1}$ is the all-one vector of length $m_1+1$, and view $A$ as a matrix over $R$. Then
\[
BB^T=\begin{pmatrix}
    AA^T & uA\mathbf{1}^T & A\mathbf{1}^T\\u\mathbf{1}A^T & u^2\mathbf{1}\cdot\mathbf{1} & u\mathbf{1}\cdot \mathbf{1}\\\mathbf{1}A^T & u\mathbf{1}\cdot \mathbf{1} & \mathbf{1}\cdot \mathbf{1}+1
\end{pmatrix}=
\begin{pmatrix}
    J_{m_1} & \mathcal{O} & \mathcal{O}\\\mathbf{0} & 0 & u\\\mathbf{0} & u & 0
\end{pmatrix}=J_{m_1}\oplus uJ_2.
\]
Define
\[
B'=\mathrm{diag}\left(
B,\underbrace{\begin{pmatrix}
1&1\\
u&0
\end{pmatrix},
\ldots,
\begin{pmatrix}
1&1\\
u&0
\end{pmatrix}
}_{(m_2-2)/2\text{ copies}}
\right).
\]
Since $B'(B')^T=J_{m_1}\oplus uJ_{m_2}$, we have $\nu(J_{m_1}\oplus uJ_{m_2})\le m_1+m_2$.

Next, we show that
\[
\nu(J_{m_1}\oplus uJ_{m_2})\ge \begin{cases}
        m_1+m_2, & \mbox{if }m_2>0,\\
        m_1+1, & \mbox{if }m_2=0.
    \end{cases}
\]
Let $\nu(J_{m_1}\oplus uJ_{m_2})=m$, and $P\in M_{n\times m}(R)$ be a matrix such that
\[
J_{m_1}\oplus uJ_{m_2}=PP^T
\]
where $n=m_1+m_2$. Let $\mathcal{C}_P$ be a code generated by the rows of $P$ with type $\{a, b\}$. Define a map $\varphi:R^{m}\to R^n$ as
\[
\varphi(\mathbf{x})=\mathbf{x}P^T
\]
for $\mathbf{x}\in R^m$. Then,
\[
\varphi(\mathcal{C}_P)=\langle PP^T\rangle=\langle J_{m_1}\oplus uJ_{m_2}\rangle=R^{m_1}\oplus (uR)^{m_2}.
\]
Since the minimal number of generators for the image of an $R$-linear map is bounded above by that of its domain, we have
\[
m_1+m_2\le a+b\le m=\nu(J_{m_1}\oplus uJ_{m_2}).
\]
Thus, $\nu(J_{m_1}\oplus uJ_{m_2})=m_1+m_2$ if $m_2>0$.

Let $m_2=0$. Then we have
\[
m_1\le \nu(J_{m_1})\le m_1+1.
\]
Suppose that there is a matrix $Q\in M_{m_1\times m_1}(R)$ such that $J_{m_1}=QQ^T$. Since $J_{m_1}$ is invertible, $Q$ is also invertible. By reducing modulo $u$, we obtain an invertible matrix $\overline{Q}$ such that $\overline{J_{m_1}}=\overline{Q}\overline{Q}^T$. Since the diagonal entries of $J_{m_1}$ are zero, every row of $\overline{Q}$ is in
\[
H=\left\{\mathbf{x}\in \mathbb{F}_q^{m_1}~|~\sum_jx_j=0\right\}
\]
whose dimension is given as $\dim H=m_1-1$. Thus $\overline{Q}$ cannot generate $\mathbb{F}_q^{m_1}$, which is a contradiction since $Q$ is invertible. Hence $\nu(J_{m_1})=m_1+1$.
\end{proof}

\end{document}
